\scanpage{76}
\mspage{75}
% E: text|upper left|par recollement des polygones d'arêtes
par recollement \uline{des} polygones \uline{d'arêtes}, \unsure{en $2n$|au $2n$}? \deleted{Cette} La première ne serait-elle pas plutôt un revêtement d'ordre deux de la seconde, ramifiée en chacun des pts du centre --- en effet, on sait \deleted{\unsure{apparemment}} \inserted{que les} antipodismes des $\Delta_D^*$ sont compatibles avec les identifications, --- par passage au quotient, cela donnera \unsure{sûrement} la carte cellulaire envisagée\unsure{\ldots} On peut alors, \deleted{qui} \unsure{pour} \deleted{\unsure{une seconde}} interpréter la nouvelle carte comme associée à un graphe \inserted{(\unsure{polygonalement})} dont les sommets sont les p.r. stables, que pour deux p.r. $D,D'$ adjacentes, on va maintenant envisager \uline{deux} flèches orientées entre eux, savoir $a_D:D\longrightarrow D'$ (\unsure{de le bord|de bords} $a$ vers $b'$) et $a_{D'}:D'\longrightarrow D$ (\unsure{de l'autre|de bords} $b$ vers $a'$), munies chacune d'une orientation --- le balayage de $D$ vers $D'$, et de $D'$ vers $D$, qu'on \unsure{désidentifie}.
% E: text|upper right|Maintnt les arcs de transition
\unsure{Maintnt|Maintt}, les arcs de transition (ancien style) pour $D$ i.e. sur $\partial\Delta_D$, sont les arcs de spécialisation \uline{et leurs antipodiques}, organisés en polygône de la façon habituelle. On considère que par l'un ou par l'autre, on va d'\unsure{une vers une} position rel. adjacentes.
% E: text|right middle|La surface Xcal obtenue
La surface $\mathcal X$ obtenue, immergée dans $X$ par une immersion, \unsure{l'image inverse} \inserted{$L$} de $X_1$ \deleted{est} donne lieu à \deleted{une} l'interprétation habituelle de la p.r. (stable) associée à une face de $(\mathcal X,K)$.
% E: text|bottom right heading|Description combinatoire
Description combinatoire
% E: display|bottom right|1 graphe Gamma sommets and arêtes orientées
\diagram{\(\begin{array}{ll}
1)\ \text{graphe }\Gamma & \left\{\begin{array}{l}
\text{sommets: p.r. stables }D\\
\text{Arêtes orientées: paires }(D,a),\ a\text{ un}\\
\qquad\text{arc de transition, sur }\underline D\\
\deleted{\text{\unsure{où}}}\qquad\text{(de spécialisation,}\ \text{ou antipodique}\\
\qquad\text{d'un tel)}
\end{array}\right.
\end{array}\)}
